\section{Experiment B: Re-rostering During Execution}
\label{sec:exp-rerostering}

% 证据：实验 4b（实验0_表示分析/reroster/：core.py、moea.py、plans.py、gen_scenarios.py、analyze.py、main_report.md、pilot_report.md）
% 注意：本实验用 κ = 0.25。与静态排班最终 κ 的关系等 κ 完整结果出来后再补一句说明。

A roster is rarely executed as planned.
Robots break down, need maintenance earlier than planned, or come back late.
The roster must then be repaired for the rest of the horizon, and the repair should change as few shifts as possible.
This is a natural case for \SDD: the original plan is a construction rule that is already known to be good,
and the task is to find the few deviations from it that the disturbance makes necessary.

\subsection{Setup}
\label{sec:exp-rerostering:setup}

\textbf{Plans and disturbances.}
For each of the 62 instances, the original plan is the roster with the fewest resources, and among those the best balance,
found by one NSGA-II run with \dec{5} (population 30, 465 evaluations).
The plan is executed up to 30\% or 60\% of its shifts.
At that point, the robot that the plan sends next is hit by one of three disturbances:
\begin{itemize}
  \item \textbf{Breakdown}: the robot is out of service for 48\,h; after the repair its duty budget is full again.
  \item \textbf{Early maintenance}: the robot must go to maintenance before its next duty, which changes the order of its tasks.
  \item \textbf{Delay}: the robot becomes ready 8\,h later than it otherwise would.
\end{itemize}
If the plan opens a new robot at that shift, the disturbance moves to the next shift.
This gives $62 \times 3 \times 2 = 372$ scenarios.
On one instance, neither simple repair of the plan (see below) is feasible in any of its six scenarios, so no budget can be set; we use the remaining 366.
The shifts before the disturbance are fixed.

\textbf{Objectives.}
A repaired roster is judged by two objectives, both minimized:
the number of changes, \ie the shifts after the disturbance whose robot differs from the planned shift with the nearest start time,
and the imbalance $\fbal$.
The resource count is a constraint: it may not exceed a budget $B$.
$B$ is the smaller resource count of two simple repairs of the plan (the plan rule below without deviations, and the plan encoded for \dec{5}),
so it is always reachable.
Rosters over budget are penalized in proportion to the excess.
In a pilot with resource count and changes as the objectives, the methods never differed in the smallest resource count (36 scenarios), so we moved it to the constraint.

\textbf{Methods.}
All methods continue decoding from the state at the disturbance, and the repair step of \dec{5} never reaches back into the fixed shifts.
\begin{itemize}
  \item \textbf{\dec{5}}: per-shift genes from a random initial population, \ie re-rostering from scratch.
  \item \textbf{\dec{5}+seed}: as \dec{5}, but 20\% of the initial population is the repaired plan and its mutants.
  This is a simplified form of the seeding of Shen and Yao \cite{shen2015dfjsp}, who fill 20\% of the population with the previous schedule and its mutants and 30\% with a repaired schedule and its mutants.
  \item \textbf{Plan as bias}: \SDD with the plan as the rule. The number of slots is set from the full horizon with $\dens = 0.25$, and the slots point only to the remaining shifts.
  The \emph{plan rule} sends the robot that the plan sends at the nearest start time if it is available,
  opens a robot of the planned type if the plan opens one there,
  and otherwise falls back to the default rule $\rul_0$.
  \item \textbf{Plan as bias + seed}: as above, and 20\% of the initial population is the genotype without deviations and its mutants.
\end{itemize}
The plan rule is fixed and not learned, so all 62 instances can be used.
Each method runs NSGA-II with population 30 and 465 evaluations, five times per scenario.

\textbf{Measures.}
The front of a run is the non-dominated set of all evaluated rosters within budget.
For each scenario, we normalize (changes, imbalance) over the solutions of all methods and compute the HV with reference point $(1.1, 1.1)$, divided by $1.1^2$.
We compare each method with \dec{5}+seed by a paired Wilcoxon test over scenarios,
and count the scenarios in which a method is significantly better or worse (Mann--Whitney, $p < 0.05$, five runs each).
We also report the smallest number of changes found and the fraction of runs that find a repair with at most three changes.

\subsection{Results}
\label{sec:exp-rerostering:results}

\begin{table}[t]
  \centering
  \caption{Re-rostering on 366 scenarios. $\Delta$: difference in mean HV to \dec{5}+seed. W/L: scenarios in which a method is significantly better/worse than \dec{5}+seed. Min.\ ch.: smallest number of changes found (mean). $\le 3$: fraction of runs that find a repair with at most three changes.}
  \label{tab:reroster}
  \footnotesize
  \setlength{\tabcolsep}{3pt}
  \begin{tabular}{@{}lcccccc@{}}
    \toprule
    Method & HV & $\Delta$ & $p$ & W/L & Min.\ ch. & $\le 3$ \\
    \midrule
    \dec{5}             & 0.339 & $-0.265$ & $7\!\times\!10^{-45}$ & 2/200  & 25.0 & 0.01 \\
    \dec{5}+seed        & 0.605 & ---      & ---                    & ---    & 9.9  & 0.36 \\
    Plan as bias        & 0.788 & $+0.184$ & $5\!\times\!10^{-26}$ & 249/44 & 7.0  & 0.28 \\
    Plan as bias + seed & 0.815 & $+0.210$ & $1\!\times\!10^{-40}$ & 238/24 & \textbf{6.4} & \textbf{0.41} \\
    \bottomrule
  \end{tabular}
\end{table}

\begin{table}[t]
  \centering
  \caption{Mean HV by disturbance type and by time of the disturbance (share of the plan already executed). Each group has 122 (type) or 183 (time) scenarios. All differences to \dec{5}+seed are significant ($p < 10^{-5}$).}
  \label{tab:reroster-groups}
  \footnotesize
  \setlength{\tabcolsep}{3pt}
  \begin{tabular}{@{}lccccc@{}}
    \toprule
    Method & Breakdown & Early maint. & Delay & 30\% & 60\% \\
    \midrule
    \dec{5}             & 0.328 & 0.329 & 0.362 & 0.335 & 0.344 \\
    \dec{5}+seed        & 0.618 & 0.645 & 0.551 & 0.546 & 0.663 \\
    Plan as bias        & 0.772 & 0.787 & 0.805 & \textbf{0.789} & 0.788 \\
    Plan as bias + seed & \textbf{0.807} & \textbf{0.832} & \textbf{0.806} & 0.782 & \textbf{0.848} \\
    \bottomrule
  \end{tabular}
\end{table}

Table~\ref{tab:reroster} gives the overall results.

\textbf{The plan works better as a bias than as a seed.}
Using the plan as the rule raises the mean HV from 0.605 to 0.788 ($+0.184$).
It is significantly better than seeding in 249 scenarios and worse in 44.
It also needs fewer changes: the smallest repair it finds changes 7.0 shifts on average, against 9.9 with seeding.
Re-rostering from scratch is far behind (0.339) and almost never finds a repair with few changes.

\textbf{The gain holds for every disturbance and time.}
In every group of Table~\ref{tab:reroster-groups}, both bias variants beat seeding, and seeding beats re-rostering from scratch.
The gain over seeding is largest for delays ($+0.254$) and smallest for early maintenance ($+0.142$).
The bias also beats seeding within each of the E, H and M families ($p \le 2 \times 10^{-4}$).

\textbf{Bias and seeding complement each other.}
Seeding finds repairs with at most three changes in 36\% of the runs, more often than the plan rule alone (28\%).
Combining the two finds repairs with at most three changes in 41\% of the runs and gives the highest mean HV (0.815),
although its HV gain over the bias alone is not significant ($p = 0.59$).
The combination helps when the disturbance comes late (60\% executed: 0.848 against 0.788, $p = 0.05$).
When it comes early, the bias alone is slightly better (0.789 against 0.782, $p = 0.003$).

\textbf{Why bias beats seeding.}
The two methods put the plan into the search in different places, and this explains the gap in changes (7.0 against 9.9).
Seeding places the plan in a few individuals of the initial population.
Under per-shift genes, a mutation of the seed changes the state, and the later genes are read in that new state;
the offspring drift away from the plan at many shifts at once (Section~\ref{sec:repr}).
With the plan as the rule, every genotype starts from the plan.
An offspring departs from it only at its active slots, and after each departure the rule returns to the plan's robot whenever that robot is available.
The whole search therefore stays near the plan, and that is where good repairs are.
